\documentclass[a4paper,12pt]{article}
\usepackage[italian]{babel}
\usepackage[utf8]{inputenc}
\usepackage[T1]{fontenc}
\usepackage{geometry}
\usepackage{hyperref}
\geometry{margin=2.5cm}

\title{\textbf{Informativa sul Trattamento dei Dati Personali}\\
ai sensi degli artt. 13 e 14 del Regolamento (UE) 2016/679 (GDPR)\\
NutriApp}
\date{\today}

\begin{document}

    \maketitle

    \section*{1. Titolare del Trattamento}

    Il Titolare del trattamento dei dati personali è \textbf{NutriApp Team}, con sede operativa in Italia.
    Per qualsiasi informazione relativa al trattamento dei dati personali è possibile contattare il Titolare al seguente indirizzo email:
    \href{mailto:privacy@nutriapp.it}{privacy@nutriapp.it}.

    Qualora venga nominato un Responsabile della Protezione dei Dati (DPO), i relativi riferimenti saranno pubblicati e comunicati agli utenti.

    \section*{2. Tipologia di Dati Trattati}

    Nell’ambito dell’utilizzo dell’applicazione NutriApp possono essere trattate le seguenti categorie di dati personali:

    \textbf{Dati identificativi e di contatto}
    nome, indirizzo email, eventuali credenziali di accesso.

    \textbf{Dati relativi alla salute (categorie particolari di dati ai sensi dell’art. 9 GDPR)}
    peso, altezza, obiettivi nutrizionali, abitudini alimentari, dati relativi all’apporto calorico e nutrizionale, eventuali informazioni fornite volontariamente dall’utente in merito a condizioni mediche o esigenze alimentari specifiche.

    \textbf{Dati di utilizzo}
    informazioni relative all’utilizzo dell’applicazione, preferenze, interazioni con le funzionalità.

    \textbf{Dati tecnici}
    indirizzo IP, identificativi del dispositivo, sistema operativo, dati di log, cookie tecnici ove applicabili.

    \section*{3. Finalità del Trattamento}

    I dati personali sono trattati per le seguenti finalità:

    \begin{itemize}
        \item Fornire i servizi di monitoraggio nutrizionale e gestione del diario alimentare;
        \item Consentire la creazione e gestione dell’account utente;
        \item Elaborare statistiche personalizzate e visualizzazioni grafiche dei progressi;
        \item Migliorare la qualità e le funzionalità dell’applicazione;
        \item Adempiere ad obblighi di legge o richieste delle autorità competenti;
        \item Gestire eventuali richieste di assistenza o supporto tecnico.
    \end{itemize}

    \section*{4. Base Giuridica del Trattamento}

    Il trattamento dei dati personali si fonda sulle seguenti basi giuridiche:

    \begin{itemize}
        \item Art. 6, par. 1, lett. b GDPR: esecuzione di un contratto o di misure precontrattuali;
        \item Art. 6, par. 1, lett. c GDPR: adempimento di obblighi legali;
        \item Art. 6, par. 1, lett. a GDPR: consenso dell’interessato per specifiche finalità;
        \item Art. 9, par. 2, lett. a GDPR: consenso esplicito per il trattamento di dati relativi alla salute.
    \end{itemize}

    Il conferimento dei dati è necessario per l’erogazione dei servizi principali dell’applicazione. Il mancato conferimento può comportare l’impossibilità di utilizzare alcune funzionalità.

    \section*{5. Modalità di Trattamento}

    Il trattamento dei dati avviene mediante strumenti informatici e telematici, secondo principi di liceità, correttezza, trasparenza, minimizzazione e limitazione della conservazione.

    Non vengono effettuati processi decisionali automatizzati che producano effetti giuridici sull’interessato ai sensi dell’art. 22 GDPR.

    \section*{6. Periodo di Conservazione dei Dati}

    I dati personali saranno conservati per il tempo strettamente necessario al conseguimento delle finalità per le quali sono stati raccolti, in conformità al principio di limitazione della conservazione di cui all’art. 5 GDPR.

    In particolare:

    \begin{itemize}
        \item I dati saranno conservati per tutta la durata del rapporto contrattuale e dell’utilizzo dell’applicazione;

        \item Una volta terminato il rapporto contrattuale o cancellato l’account, i dati personali potranno essere conservati per un periodo ulteriore qualora ciò sia necessario per:
        \begin{itemize}
            \item adempiere ad obblighi di legge (ad esempio obblighi fiscali o contabili);
            \item tutelare i diritti del Titolare, anche in sede giudiziaria;
            \item gestire eventuali reclami o controversie.
        \end{itemize}

        \item I dati relativi alla salute saranno conservati esclusivamente per il tempo necessario alla fornitura del servizio richiesto dall’utente e saranno cancellati o anonimizzati in caso di cessazione del rapporto, salvo sussista un obbligo di legge o altra base giuridica che ne consenta la conservazione.
    \end{itemize}

    Al termine dei periodi sopra indicati, i dati saranno cancellati in modo sicuro o resi anonimi in maniera irreversibile.

    \section*{7. Comunicazione e Trasferimento dei Dati}

    I dati personali potranno essere comunicati a:

    \begin{itemize}
        \item Fornitori di servizi informatici e hosting;
        \item Fornitori di infrastrutture cloud;
        \item Consulenti legali o fiscali, ove necessario.
    \end{itemize}

    Tali soggetti operano in qualità di Responsabili del trattamento ai sensi dell’art. 28 GDPR.

    Qualora i dati vengano trasferiti verso Paesi terzi al di fuori dello Spazio Economico Europeo, il trasferimento avverrà nel rispetto degli artt. 44 e seguenti del GDPR, mediante l’adozione di adeguate garanzie (ad esempio Clausole Contrattuali Standard).

    \section*{8. Diritti dell’Interessato}

    L’interessato può esercitare in qualsiasi momento i diritti previsti dagli artt. 15–22 GDPR, tra cui:

    \begin{itemize}
        \item Diritto di accesso ai dati personali;
        \item Diritto di rettifica;
        \item Diritto alla cancellazione (diritto all’oblio);
        \item Diritto alla limitazione del trattamento;
        \item Diritto alla portabilità dei dati;
        \item Diritto di opposizione al trattamento;
        \item Diritto di revoca del consenso in qualsiasi momento.
    \end{itemize}

    Le richieste possono essere inviate all’indirizzo email \href{mailto:privacy@nutriapp.it}{privacy@nutriapp.it}.

    È inoltre possibile proporre reclamo all’Autorità Garante per la Protezione dei Dati Personali.

    \section*{9. Sicurezza dei Dati}

    NutriApp adotta misure tecniche e organizzative adeguate ai sensi dell’art. 32 GDPR, tra cui:

    \begin{itemize}
        \item Autenticazione sicura degli utenti;
        \item Accesso ai dati limitato al personale autorizzato;
        \item Backup periodici e sistemi di protezione contro accessi non autorizzati.
    \end{itemize}

    \section*{10. Minori}

    L’applicazione non è destinata a minori di anni 14.

    Ai sensi dell’art. 8 del GDPR e dell’art. 2-quinquies del D.lgs. 196/2003 (come modificato dal D.lgs. 101/2018), il trattamento dei dati personali basato sul consenso in relazione all’offerta diretta di servizi della società dell’informazione è lecito solo se il minore ha compiuto almeno 14 anni.

    Qualora l’utente abbia un’età inferiore ai 14 anni, il trattamento dei dati personali è lecito soltanto se il consenso è prestato o autorizzato dall’esercente la responsabilità genitoriale.

    Il Titolare potrà adottare misure ragionevoli per verificare l’età dichiarata dall’utente e, ove necessario, la sussistenza del consenso del genitore o tutore.

    Qualora venga accertato che siano stati raccolti dati personali di un minore di 14 anni in assenza del consenso richiesto, tali dati saranno tempestivamente cancellati.

    \section*{11. Modifiche alla Presente Informativa}

    La presente informativa potrà essere aggiornata nel tempo.
    Eventuali modifiche sostanziali saranno comunicate agli utenti tramite notifica in-app o via email.

    \vspace{1cm}

    \noindent
    Ultimo aggiornamento: \today

\end{document}
